% Canonical LaTeX source; imported and revised from the legacy Markdown.
\chapter{Teoremi del calcolo differenziale}
\label{chap:11}

\begin{obiettivocapitolo}{scegliere il teorema corretto e verificarne le ipotesi}
\prioritaalta. Nei problemi di teoremi, la verifica delle ipotesi è parte della soluzione, non un dettaglio formale.
\end{obiettivocapitolo}
\begin{sceltametodo}{quale teorema usare}
\begin{tabularx}{\textwidth}{@{}p{.25\textwidth}p{.31\textwidth}X@{}}
Fermat & estremo locale interno e derivabilità & conclude $f'(x_0)=0$; è condizione necessaria, non sufficiente.\\
Rolle & continuità, derivabilità, $f(a)=f(b)$ & produce un punto con tangente orizzontale.\\
Lagrange & continuità e derivabilità & collega variazione media e derivata; utile per monotonia, stime e unicità.\\
Cauchy & due funzioni e rapporto di incrementi & generalizza Lagrange; base teorica di de l'Hôpital.\\
De l'Hôpital & forma $0/0$ o $\infty/\infty$ & deriva numeratore e denominatore, non il quoziente; il nuovo limite deve esistere.
\end{tabularx}
\end{sceltametodo}
\begin{proceduraesame}{applicazione di un teorema}
\begin{enumerate}
\item Scrivi l'intervallo e identifica estremi e interno.
\item Verifica continuità sul chiuso e derivabilità sull'aperto, oltre alle ipotesi specifiche.
\item Cita il teorema e scrivi la conclusione nella forma esatta.
\item Se la traccia chiede il punto $c$, risolvi l'equazione prodotta dal teorema e seleziona le soluzioni in $(a,b)$.
\end{enumerate}
\end{proceduraesame}
\begin{errorefrequente}{conversi falsi}
$f'(x_0)=0$ non implica che $x_0$ sia un estremo; la continuità non implica derivabilità; il soddisfacimento della conclusione non dimostra le ipotesi del teorema.
\end{errorefrequente}

Nei teoremi seguenti le righe prima della freccia sono le \textbf{ipotesi} e ciò che segue la freccia è la \textbf{conclusione}. Il converso non è valido salvo quando è espressamente dichiarata un'equivalenza.

\section{Teorema di Fermat}\label{teorema-di-fermat}

\[
\displaystyle
\left.
\begin{array}{l}
x_0\in\operatorname{int}D,\\
x_0\text{ estremo locale di }f,\\
f\text{ derivabile in }x_0
\end{array}
\right\}
\quad\Longrightarrow\quad
f'(x_0)=0.
\]

Condizione necessaria, non sufficiente:

\[
\displaystyle
f(x)=x^3,
\qquad
f'(0)=0,
\qquad
0\text{ non è un estremo}.
\]

Caso di frontiera:

\[
\displaystyle
f(x)=x\text{ su }[0,1],
\qquad
\max f=f(1),
\qquad
f'_-(1)=1.
\]

Approfondimento: \href{https://ingegnerismo.it/matematica/teorema-di-fermat/}{teorema di Fermat}.

\section{Teorema di Rolle}\label{teorema-di-rolle}

Per \(a<b\):

\[
\displaystyle
\left.
\begin{array}{l}
f\in C([a,b]),\\
f\text{ derivabile in }(a,b),\\
f(a)=f(b)
\end{array}
\right\}
\quad\Longrightarrow\quad
\exists c\in(a,b):\ f'(c)=0.
\]

Approfondimento: \href{https://ingegnerismo.it/matematica/teorema-di-rolle/}{teorema di Rolle}.

\section{Teorema di Lagrange}\label{teorema-di-lagrange}

Per \(a<b\):

\[
\displaystyle
\left.
\begin{array}{l}
f\in C([a,b]),\\
f\text{ derivabile in }(a,b)
\end{array}
\right\}
\quad\Longrightarrow\quad
\exists c\in(a,b):
f'(c)=\frac{f(b)-f(a)}{b-a}.
\]

Se \(f\) è derivabile su un intervallo \(I\), valgono le conseguenze:

\[
\begin{aligned}
f'\ge0\text{ su }I&\Longrightarrow f\text{ non decrescente su }I,\\
f'\le0\text{ su }I&\Longrightarrow f\text{ non crescente su }I,\\
f'=0\text{ su }I&\Longrightarrow f\text{ costante su }I,\\
f'>0\text{ su }I&\Longrightarrow f\text{ strettamente crescente su }I,\\
f'<0\text{ su }I&\Longrightarrow f\text{ strettamente decrescente su }I.
\end{aligned}
\]

Se \(|f'|\le M\) su \(I\):

\[
\displaystyle
|f(y)-f(x)|\le M|y-x|
\qquad\forall x,y\in I.
\]

Approfondimento: \href{https://ingegnerismo.it/matematica/teorema-di-lagrange/}{teorema di Lagrange}.

Esercizi svolti: \href{https://ingegnerismo.it/matematica/fermat-rolle-lagrange-esercizi/}{Fermat, Rolle e Lagrange}.

\section{Teorema di Cauchy}\label{teorema-di-cauchy}

Per \(a<b\):

\[
\displaystyle
\left.
\begin{array}{l}
f,g\in C([a,b]),\\
f,g\text{ derivabili in }(a,b)
\end{array}
\right\}
\quad\Longrightarrow\quad
\exists c\in(a,b):
\]

\[
\displaystyle
f'(c)[g(b)-g(a)]
=
g'(c)[f(b)-f(a)].
\]

Se inoltre \(g'(x)\ne0\) per ogni \(x\in(a,b)\):

\[
\displaystyle
\exists c\in(a,b):
\frac{f'(c)}{g'(c)}
=
\frac{f(b)-f(a)}{g(b)-g(a)}.
\]

Caso particolare \(g(x)=x\): teorema di Lagrange.

Approfondimento: \href{https://ingegnerismo.it/matematica/teorema-di-cauchy/}{teorema di Cauchy}.

Esercizi svolti: \href{https://ingegnerismo.it/matematica/teoremi-calcolo-differenziale-esercizi/}{teorema di Cauchy}.

\section{Teorema di Darboux per le derivate}\label{teorema-di-darboux-per-le-derivate}

Sia \(f\) derivabile su un intervallo \(I\); per \(x_1<x_2\) in \(I\) e

\[
\min\{f'(x_1),f'(x_2)\}<\lambda<\max\{f'(x_1),f'(x_2)\},
\]

vale

\[
\exists c\in(x_1,x_2):\ f'(c)=\lambda.
\]

\[
f'(I)\text{ è un intervallo}.
\]

Approfondimento: \href{https://ingegnerismo.it/matematica/teorema-di-darboux/}{teorema di Darboux}.

Esercizi svolti: \href{https://ingegnerismo.it/matematica/teorema-di-darboux-esercizi/}{teorema di Darboux per le derivate}.

\section{Teorema di de l'Hôpital}\label{teorema-di-de-lhuxf4pital}

Sia \(a\in\overline{\mathbb R}\) e sia \(I\) un intervallo puntato, eventualmente unilatero, che tende ad \(a\):

\[
\displaystyle
f,g\text{ derivabili in }I,
\qquad
g(x)\ne0,
\qquad
g'(x)\ne0
\quad\text{in }I.
\]

Forme ammesse:

\[
\displaystyle
f(x)\to0,
\quad g(x)\to0,
\]

oppure

\[
\displaystyle
|f(x)|\to+\infty,
\quad |g(x)|\to+\infty.
\]

Se

\[
\displaystyle
\lim_{x\to a}\frac{f'(x)}{g'(x)}=L
\in\overline{\mathbb R},
\]

allora

\[
\displaystyle
\lim_{x\to a}\frac{f(x)}{g(x)}=L.
\]

Riduzioni standard:

\begin{longtable}[]{@{}ll@{}}
\toprule\noalign{}
Forma & Trasformazione \\
\midrule\noalign{}
\endhead
\bottomrule\noalign{}
\endlastfoot
\(0\cdot\infty\) & \(u v=\dfrac{u}{1/v}=\dfrac{v}{1/u}\) \\
\(\infty-\infty\) & denominatore comune oppure razionalizzazione \\
\(0^0\), \(1^\infty\), \(\infty^0\) & \(u^v=\exp(v\ln u)\), con \(u>0\) \\
\end{longtable}

Approfondimento: \href{https://ingegnerismo.it/matematica/teorema-di-de-l-hopital/}{teorema di de l'Hôpital}.
